Next we consider the slice rank of a linear subspace of tensors.  Given a linear subspace $\mathsf{W} \subseteq \mathbb{K}^{n_1} \otimes_{\mathbb{K}}  \cdots \otimes_{\mathbb{K}}   \mathbb{K}^{n_d}$ over a field $\mathbb{K}$,  the \emph{slice rank of $\mathsf{W}$} is defined as 
\begin{equation}\label{eq:SRW}
\sr_{\mathbb{K}}(\mathsf{W}) \coloneqq \min \left\lbrace \sum_{i=1}^{d} \codim( \mathsf{U}_{i}):  \langle \mathsf{W}, \mathsf{U}_{1} \otimes_{\mathbb{K}}  \cdots \otimes_{\mathbb{K}}  \mathsf{U}_{d}\rangle=0,\; \mathsf{U}_i \subseteq (\mathbb{K}^{n_i})^{\ast},\;1 \le i  \le d \right\rbrace.
\end{equation}
Here $\langle \mathsf{W}, \mathsf{U}_{1}\otimes \cdots\otimes  \mathsf{U}_{d}\rangle = 0$ means that $\langle T,  L \rangle = 0$ for each pair $(T, L) \in \mathsf{W} \times \left( \mathsf{U}_{1}\otimes \cdots\otimes \mathsf{U}_{d} \right)$.  In the sequel,  we denote 
\begin{equation}\label{eq:Vi}
\mathsf{V}_i \coloneqq \mathbb{K}^{n_1} \otimes_{\mathbb{K}}  \cdots \otimes_{\mathbb{K}}  \mathbb{K}^{n_{i-1}} \otimes_{\mathbb{K}}  \mathbb{K}^{n_{i+1}} \otimes_{\mathbb{K}}  \cdots \otimes_{\mathbb{K}}  \mathbb{K}^{n_d},
\end{equation}
for each $1 \le i \le d$.
\begin{lemma}[Characterization of $\sr_{\mathbb{K}}(\mathsf{W})$]\label{lem4-8}
Let $\mathbb{K}$ be a field.  For a linear subspace $\mathsf{W} \subseteq \mathbb{K}^{n_1} \otimes_{\mathbb{K}}  \cdots \otimes_{\mathbb{K}}  \mathbb{K}^{n_d}$,  $\sr_{\mathbb{K}}(\mathsf{W})$ is the minimal positive integer $r$ to ensure the existence of $r_i \in \mathbb{N}$ and $a_{i1},\dots,  a_{ir_i} \in \mathbb{K}^{n_i}$ for $1 \le i \le d$,  such that $r = \sum_{i=1}^d r_i$ and every $T \in \mathsf{W}$ admits a decomposition
\begin{equation}\label{lem4-8:eq1}
T = \sum_{i=1}^d  \sum_{j=1}^{r_i} a_{ij} \otimes T_{ij}
\end{equation}
for some $T_{i1},\dots,  T_{ir_i} \in \mathsf{V}_i$,  $1 \le i \le d$.  
\end{lemma}
\begin{proof}
We denote by $s$ the minimal integer such that \eqref{lem4-8:eq1} holds.  We prove that $s = t \coloneqq \sr_{\mathbb{K}}(\mathsf{W})$.  If $\langle \mathsf{W}, \mathsf{U}_{1} \otimes_{\mathbb{K}}  \cdots \otimes_{\mathbb{K}}  \mathsf{U}_{d}\rangle=0$,  then $\mathsf{W} \subseteq \sum_{i=1}^d \mathsf{U}_{i}^{\perp} \otimes \mathsf{V}_{i}$ where $\mathsf{U}_{i}^{\perp} \coloneqq ( (\mathbb{K}^{n_i})^\ast/\mathsf{U}_i)^\ast$.  Hence any $T\in \mathsf{W}$ admits a decomposition \eqref{lem4-8:eq1} if we choose $r_i = \dim \mathsf{U}_i^{\perp}$ and let $\{a_{i1},\dots,  a_{ir_i}\}$ be a basis of $\mathsf{U}_i^{\perp}$ for $1 \le i \le d$.  Thus we have $t = \sum_{i=1}^d \dim \mathsf{U}_i^\perp = \sum_{i=1}^d r_i \ge s$.

Conversely,  suppose that $r_i$'s and $a_{ij}$'s are given so that $s = \sum_{i=1}^d r_i$ and \eqref{lem4-8:eq1} holds for each $T \in \mathsf{W}$.  We define 
\[
\mathsf{U}_i \coloneqq \left( \spa  \{ a_{i1},\dots,  a_{ir_i} \} \right)^\perp  \coloneqq \left(\mathbb{K}^{n_i}/\spa  \{ a_{i1},\dots,  a_{ir_i} \} \right)^\ast \subseteq (\mathbb{K}^{n_i})^\ast,\quad 1 \le i \le d.
\]
Then clearly we have $\langle \mathsf{W}, \mathsf{U}_{1} \otimes_{\mathbb{K}}  \cdots \otimes_{\mathbb{K}}   \mathsf{U}_{d}\rangle = 0$ and this implies $s = \sum_{i=1}^d r_i \ge t$.
\end{proof}
\begin{lemma}\label{lem4-9}
For any field $\mathbb{K}$ and linear subspace $\mathsf{W} \subseteq \mathbb{K}^{n_1} \otimes_{\mathbb{K}}  \cdots \otimes_{\mathbb{K}}  \mathbb{K}^{n_d}$,  there exist $n_{d+1},n_{d+2} \in \mathbb{N}$ and $T_{\mathsf{W}}\in \mathbb{K}^{n_1} \otimes_{\mathbb{K}}   \cdots \otimes_{\mathbb{K}}  \mathbb{K}^{n_{d+2}}$ such that $\sr_{\mathbb{K}}(\mathsf{W}) = \sr_{\mathbb{K}}(T_{\mathsf{W}})$. 
\end{lemma}
\begin{proof} 
Assume $\dim \mathsf{W} = n$.  Let $T_1,\dots,  T_n$ be a basis of $\mathsf{W}$.  We denote $m \coloneqq \max \left\lbrace n_i: 1 \le i \le d \right\rbrace+1$.  Suppose that $\{e_{ij}: 1 \le i \le m,\; 1 \le j \le n \}$ is a basis of $\mathbb{K}^{mn}$.  We consider
\[
T_{\mathsf{W}} \coloneqq \sum_{j=1}^{n} T_{j}\otimes  \left( \sum_{i=1}^{m}  e_{ij} \otimes e_{ij} \right) \in \mathbb{K}^{n_1} \otimes_{\mathbb{K}}   \cdots \otimes_{\mathbb{K}}  \mathbb{K}^{n_d} \otimes_{\mathbb{K}}  \mathbb{K}^{mn} \otimes_{\mathbb{K}}  \mathbb{K}^{mn}. 
\]
By Lemma~\ref{lem4-8},  it is clear that $\sr_{\mathbb{K}}(T_{\mathsf{W}}) \le \sr_{\mathbb{K}}(\mathsf{W})$.  Thus, it suffices to prove $\sr_{\mathbb{K}}(T_{\mathsf{W}}) \ge \sr_{\mathbb{K}}(\mathsf{W})$. 

Let $T_{\mathsf{W}}=\sum_{i=1}^{d+2}\sum_{j=1}^{r_i} a_{ij}\otimes S_{ij}$ be a slice rank decomposition of $T_{\mathsf{W}}$,  where $a_{ij} \in \mathbb{K}^{n_i}$ and $S_{ij} \in \mathsf{V}_i$ for $1 \le i \le d+2$ and $1\le j \le r_i$.  Here $\mathsf{V}_i$ is defined as in \eqref{eq:Vi}.  We denote 
\[
\mathsf{W}' \coloneqq \sum_{i=1}^{d }\sum_{j=1}^{r_i}  \left\langle a_{ij}\otimes \mathsf{V}_i,  (\mathbb{K}^{mn})^\ast \otimes_{\mathbb{K}}  (\mathbb{K}^{mn})^\ast \right\rangle \subseteq \mathbb{K}^{n_1} \otimes_{\mathbb{K}}  \cdots \otimes_{\mathbb{K}}  \mathbb{K}^{n_d}.
\]
%
%
%
%
We observe that each element of $\mathsf{W}'$ is a linear combination of $a_{ij} \otimes R_{ij}$ for some 
\[
R_{ij} \in \langle \mathsf{V}_{i},  (\mathbb{K}^{mn})^\ast \otimes_{\mathbb{K}}  (\mathbb{K}^{mn})^\ast \rangle, \; 1\le i \le d,\; 1 \le j \le r_i.
\]
Without loss of generality,  we may assume $T_{1},\cdots, T_{l}\in \mathsf{W}'$ and $T_{l+1},\cdots,T_{n}\notin \mathsf{W}'$ for some $0 \le l \le n$. 
In the following,  we regard $a_{ij}$'s (resp.  $T_k$'s) as linear (resp.  multilinear) polynomials.  Since $T_{l+1}, \dots,  T_{n} \not\in \mathsf{W}'$,  we conclude that $T_{l+1}\cdots T_{n}  \notin J$, where $J$ is the prime ideal generated by  linear forms $a_{ij},  1\le i \le d,  1 \le j \le r_i$.  By Hilbert Nullstellensatz,  there exist $v_i\in \overline{\mathbb{K}}^{n_i}$,  $1\le i\le d$, such that 
\[
a_{ij}(v_{i})=0,\; T_{s}(v_{1},\cdots,v_{d})\not=0,\;1\le j\le r_{i},\; l+1\le s\le d.
\]
Therefore,  we obtain 
\begin{align}
T_{\mathsf{W}}(v_{1},\cdots,v_{d})&=\sum_{i=1}^{d+2}\sum_{j=1}^{r_{i}} (a_{ij} \otimes S_{ij})( v_1,\dots, v_{d} ) \nonumber \\
&= \sum_{j=1}^{r_{d+1}}a_{d+1,j} \otimes S_{d+1,j}(v_1,\dots, v_{d})
+ 
\sum_{j=1}^{r_{d+2}}a_{d+2,j} \otimes S_{d+2,j}(v_1,\dots, v_{d})
\label{lem4-9:eq1}
\end{align}
On the other hand,  by the construction of $T_{\mathsf{W}}$,  we also have 
\begin{equation}\label{lem4-9:eq2}
T_{\mathsf{W}}(v_{1},\cdots,v_{d})=
\sum_{s=1}^{n} T_{s} (v_1,\dots,  v_d)  \left( \sum_{i=1}^{m}  e_{is}\otimes e_{is} \right) = \sum_{s=l+1}^{n} T_{s} (v_1,\dots,  v_d)  \left( \sum_{i=1}^{m}  e_{is}\otimes e_{is} \right).
\end{equation}
Both \eqref{lem4-9:eq1} and \eqref{lem4-9:eq2} are decompositions of the matrix $T_{\mathsf{W}}(v_{1},\cdots,v_{d}) \in \overline{\mathbb{K}}^{mn \times mn}$ into the sum of rank one matrices.  Moreover,  \eqref{lem4-9:eq2} is clearly a rank decomposition. This implies that $(n-l)m \le r_{d+1} + r_{d+2}$ which implies $l=n$ since $r_{d+1}+r_{d+2}\le\sr_{\mathbb{K}}(T_{\mathsf{W}}) <  m$ by construction. Hence all $T_{i}$ are in $\mathsf{W}'$,  and this implies $\operatorname{SR}_{\mathbb{K}} (\mathsf{W}) \le \operatorname{SR}_{\mathbb{K}} (\mathsf{W}') \le \operatorname{SR}_{\mathbb{K}} ( T_\mathsf{W})$.
\end{proof}
